\documentclass[tikz,border=6pt]{standalone}
\usepackage{mathdiagrams}
% isosceles_triangle_xyz_42 — XY = YZ, ∠XYZ = 42°, base angles = x
% Contract: no white boxes; no prose captions (formulas only).
% CRITICAL: TikZ `angle=A--B--C` draws counterclockwise from BA to BC.
%   At X, `angle=Y--X--Z` runs 50° → 0° the LONG way (310° reflex).
%   Use `angle=Z--X--Y` so it sweeps the small interior angle.
\begin{document}
\begin{tikzpicture}[md]
  \coordinate (Y) at (0,2.6);
  \coordinate (X) at (-2.2,0);
  \coordinate (Z) at (2.2,0);
  \draw[md shape] (X) -- (Y) -- (Z) -- cycle;

  % equal-side ticks on XY and YZ
  \draw[md thin] ($(X)!0.5!(Y)+(0.14,0.10)$) -- ($(X)!0.5!(Y)+(-0.10,0.14)$);
  \draw[md thin] ($(Y)!0.5!(Z)+(0.10,0.14)$) -- ($(Y)!0.5!(Z)+(-0.14,0.10)$);

  % vertices — plain, no boxes
  \node[font=\sffamily\small, text=mdInk] at (Y) [above] {$Y$};
  \node[font=\sffamily\small, text=mdInk] at (X) [below left] {$X$};
  \node[font=\sffamily\small, text=mdInk] at (Z) [below right] {$Z$};

  % angles — ray order chosen so the SMALL interior angle is drawn
  \pic[md angle arc, angle radius=0.55cm] {angle=X--Y--Z};
  \pic[md angle arc, angle radius=0.55cm] {angle=Z--X--Y};
  \pic[md angle arc, angle radius=0.55cm] {angle=Y--Z--X};

  % angle values — inside each angle, plain text
  \node[font=\sffamily\small, text=mdInk] at (0,1.85) {$42^\circ$};
  \node[font=\sffamily\small, text=mdInk] at (-1.45,0.55) {$x$};
  \node[font=\sffamily\small, text=mdInk] at (1.45,0.55) {$x$};

  % formulas only — no prose
  \node[font=\sffamily\small, text=mdInk] at (0,-1.15)
    {$XY = YZ \qquad 2x + 42^\circ = 180^\circ$};
\end{tikzpicture}
\end{document}
